%\documentclass[12pt,oneside]{amsart}
\documentclass{scrartcl}

%\documentclass{article}

\usepackage{amsthm,amsmath,amssymb}
\usepackage[numbers]{natbib}

%\usepackage[colorlinks]{hyperref}
\usepackage{hyperref}
\hypersetup{
    colorlinks,%
    citecolor=black,%
    filecolor=black,%
    linkcolor=black,%
    urlcolor=black
}
\usepackage{hypernat}
\usepackage[top=2.5cm, bottom=2.5cm, left=2.8cm,
right=2.8cm]{geometry}
\usepackage{graphicx}


%\setlength{\parskip}{0pt}
%\setlength{\parsep}{0pt}
%\setlength{\headsep}{0pt}
%\setlength{\topskip}{0pt}
%\setlength{\topmargin}{0pt}
%\setlength{\topsep}{0pt}
%\setlength{\partopsep}{0pt}

\linespread{1}

%\usepackage{marvosym}

%\textwidth = 6.5 in \textheight = 9.2 in \oddsidemargin = 0.0 in
%\evensidemargin = 0.0 in \topmargin = -0.0 in \headheight = 0.0 in
%\headsep = 0.0in \parskip = 0.0in \parindent = 0.0in
%\def\baselinestretch{0.871}


\newtheorem{theorem}{Theorem}[section]
\newtheorem{proposition}[theorem]{Proposition}
\newtheorem{problem}[theorem]{Problem}
\newtheorem{fact}[theorem]{Fact}
\newtheorem{lemma}[theorem]{Lemma}
\newtheorem{defn}[theorem]{Definition}
\newtheorem{corollary}[theorem]{Corollary}
\newtheorem{remark}{Remark}[section]
\newtheorem{example}[theorem]{Example}

\newcommand{\E}{{\mathbb E}}
\newcommand{\se}{{\mathcal{E}}}
\newcommand{\R}{{\mathbb R}}
\renewcommand{\P}{{\mathbb P}}
\newcommand{\ac}{{\mathcal{A}}}
\newcommand{\A}{{\mathcal{A}}}
\newcommand{\B}{{\mathcal{B}}}
\newcommand{\C}{{\mathcal{C}}}
\newcommand{\M}{{\mathcal{M}}}
\newcommand{\Mf}{{\mathfrak{M}}}
\newcommand{\T}{{\mathcal{T}}}
\newcommand{\Z}{{\mathcal{Z}}}
\newcommand{\K}{{\mathcal{K}}}
\newcommand{\lc}{{\mathcal{L}}}
\newcommand{\Ps}{{\mathcal{P}}}
\newcommand{\G}{{\mathcal{G}}}
\newcommand{\pa}{{\mathring{p}}}
\newcommand{\F}{{\cal F}}
\newcommand{\samp}{{\mathcal{X}}}
\newcommand{\X}{{\mathcal{X}}}
\newcommand{\hi}{{\hat{\Phi}}}
\newcommand{\data}{{\text{data}}}
\newcommand{\relu}{{\text{ReLU}}}

\newcommand{\ridge}{{\hat{\beta}_{\text{ridge}}(\lambda)}}
\newcommand{\lasso}{{\hat{\beta}_{\text{lasso}}(\lambda)}}
\newcommand{\ridgemu}{{\hat{\mu}_t^{\text{ridge}}(\lambda)}}
\newcommand{\lassomu}{{\hat{\mu}_t^{\text{lasso}}(\lambda)}}


\newcounter{rcnt}[section]
\renewcommand{\thercnt}{(\roman{rcnt})}

\def\qt#1{\qquad\text{#1}}

\def\argmin{\mathop{\rm argmin}}
\def\argmax{\mathop{\rm argmax}}


\setlength{\parskip}{1.4 \medskipamount}
%\sloppy
%\linespread{1.3}

\begin{document}

\title{STAT 238 - Bayesian Statistics  \\
  Lecture Thirty One} 
\subtitle{Spring 2026, UC Berkeley}

\author{Aditya Guntuboyina}


\date{15 April 2026}

\maketitle

We will study the MALA and HMC algorithms. MALA stands for Metropolis
Adjusted Langevin Algorithm. HMC stands for Hamiltonian Monte
Carlo. Before we discuss MALA, let us first recall the basics of
markov chains. 

\section{Recap: Markov Chains, detailed balance and the
  Metropolis-Hastings Algorithm}

Let $S$ be a state space. A sequence of random variables $\Theta^{(0)},
\Theta^{(1)}, \dots$ taking values in $S$ is a 
Markov Chain on $S$ provided the conditional distribution of  
$\Theta^{(t+1)}$ given $\Theta^{(t)} = \theta^{(t)}, \Theta^{(t-1)} =
\theta^{(t-1)}, \dots, \Theta^{(0)} = \theta^{(0)}$ only depends on
$\theta^{(t)}$. Further, if this conditional distribution is the same
for all  $t$, then the Markov chain is said to be
time-homogeneous. We will only deal with time-homogeneous Markov
Chains in this course.

A (time-homogeneous) Markov Chain is described via the conditional
distribution of $\Theta^{(t+1)}$ given $\Theta^{(t)} = x$, which is
described by a probability transition kernel $P(x, y)$. If $S$ is
discrete, $P(x, y)$ is the conditional probability of $\Theta^{(t+1)}
= y$ given $\Theta^{(t)} = x$. If $S$ is an open subset of $\R^d$,
$P(x, y)$ represents the conditional density of $\Theta^{(t+1)}$ given
$\Theta^{(t)} = x$.

The Markov chain given by $P(x, y)$ satisfies \textbf{detailed
  balance} with respect to a probability $\pi(x)$ on $S$ if the
following condition is satisfied:
\begin{align}\label{det_bal}
  \pi(y) P(y, x) = \pi(x) P(x, y) ~~\text{for all $x, y \in S$}. 
\end{align}
This detailed balance condition implies that $\pi$ is stationary for
the chain $P$ in the sense that $\Theta^{(t)} \sim \pi$ implies
$\Theta^{(t+1)} \sim \pi$.

Stationarity along with the additional condition of irreducibility
(which means that for every pair of points $x$ and $y$ in the state
space, there is a positive probability that the chain goes from $x$ to
$y$ in finite time) imply ergodicity which means:
\begin{align*}
  \lim_{N \rightarrow \infty} \frac{1}{N} \sum_{t=1}^N g(\Theta^{(t)})
  = \int g(\theta) d\pi(\theta)
\end{align*}
almost surely.

Metropolis-Hastings is the most widely used method for constructing
markov chains satisfying detailed balance with respect to $\pi$. This
method starts with an arbitrary transition kernel $Q(x, y)$ that may
not even have anything to do with $\pi$ (in particular, it is not at
all necessary that $Q$ satisfies detailed balance with respect to
$\pi$). Metropolis-Hastings uses $Q(\cdot, \cdot)$ and $\pi(\cdot)$ to
  construct a new Markov chain $\{\Theta^{(t)}\}$ which moves as
  follows. Given $\Theta^{(t)} = x$, first use $Q(x, \cdot)$ to
  generate $y$, then take $\Theta^{(t+1)}$ to be:
  \[
\Theta^{(t+1)} =
\begin{cases}
y, & \text{with probability } \min\!\left(1, \frac{\pi(y)\,Q(y,x)}{\pi(x)\,Q(x,y)}\right), \\[1em]
x, & \text{with probability } 1 - \min\!\left(1, \frac{\pi(y)\,Q(y,x)}{\pi(x)\,Q(x,y)}\right).
\end{cases}
\]
The resulting Metropolis-Hastings chain  satisfies detailed
balance with respect to $\pi$, as we saw previously in Lecture 25.

\section{Random Walk Metropolis (RWM)}

The Random Walk Metropolis (RWM) algorithm is arguably the simplest
example of a markov chain which satisfies detailed balance with
respect to $\pi$. It uses the transition kernel $Q(x, y)$ which
generates proposals via:
\begin{align}\label{rwm_y}
  y = x + \sigma z ~~ \text{where $z \sim N(0, I_d)$}. 
\end{align}
In other words, the proposal transition kernel is given by
\begin{align}\label{rwm_prop}
  Q(x, y) = \left(\frac{1}{\sqrt{2 \pi} \sigma} \right)^d \exp
  \left(-\frac{\|x - y\|^2}{2\sigma^2} \right) . 
\end{align}
Clearly $Q(x, y) = Q(y, x)$ so the resulting acceptance probability
(in the Metropolis-Hastings algorithm) is simply
\begin{align*}
    \min \left(1, \frac{\pi(y)}{\pi(x)} \right). 
\end{align*}
The choice of $\sigma$ (in \eqref{rwm_prop}) is crucial to the
practical performance of the algorithm. If $\sigma$ is not small, then
$Q(x, \cdot)$ may result in proposals $y$ that are far from $x$. If
$\pi(\cdot)$ is highly concentrated (as is 
typically the case with posteriors), then $\pi(y)$ might be much
smaller than $\pi(x)$ leading to the chain rejecting often and getting
stuck. On the other hand if $\sigma$ is too small, then the chain
might take a long time to explore the full distribution $\pi(\cdot)$.

\section{Metropolis Adjusted Langevin Algorithm (MALA)}

In MALA, the rwm proposal \eqref{rwm_y} is modified with an expression
that depends on $\pi$. To find this formula, let us first consider the
following proposal: 
\begin{align}\label{rwm_g}
  y = x + g(x) + \sigma z ~~\text{with $z \sim N(0, I_d)$}
\end{align}
for some function $g(x)$ depending on $\pi(x)$. Clearly \eqref{rwm_y}
can be seen as a special case of \eqref{rwm_g} corresponding to $g(x)
= 0$. The proposal transition kernel corresponding to \eqref{rwm_g} is
given by: 
\begin{align*}
  Q(x, y) = \left(\frac{1}{\sqrt{2 \pi} \sigma} \right)^d
 \exp \left(-\frac{1}{2 \sigma^2} \|y - x - g(x)
  \|^2 \right).
\end{align*}
Note that $Q(x, y)$ is no longer equal to $Q(y, x)$. The key question
is how to choose $g(x)$. It is natural to choose $g(x)$ so that the
detailed balance condition holds with respect to $\pi$: 
\begin{align*}
  \pi(x) Q(x, y) \approx \pi(y) Q(y, x) 
\end{align*}
at least approximately when $y$ is close to $x$. Check that
\begin{align*}
  \frac{Q(y, x)}{Q(x, y)} &= \exp \left(-
                                              \frac{1}{\sigma^2}\left<y
                                              - 
                                              x, g(x) + g(y) \right> +
                                              \frac{1}{2 \sigma^2}  
  \left(\|g(x)\|^2 - \|g(y)\|^2 \right) \right). 
\end{align*}
As a result:
\begin{align*}
&  \frac{\pi(y)Q(y, x)}{\pi(x) Q(x, y)} \\ &= \exp
                                                           \left(\log
                                                           \pi(y) -
                                                           \log \pi(x)
                                                           \right) \exp \left(-
                                              \frac{1}{\sigma^2}\left<y
                                              - 
                                              x, g(x) + g(y) \right> +
                                              \frac{1}{2 \sigma^2}  
  \left(\|g(x)\|^2 - \|g(y)\|^2 \right) \right). 
\end{align*}
For $y$ close to $x$, we use the approximations: 
\begin{align*}
&  \log \pi(y) - \log \pi(x) \approx \left<y - x, \nabla \log \pi(x)
  \right> \\
&   \frac{1}{\sigma^2}\left<y
                                              - 
                                              x, g(x) + g(y) \right>
  \approx \frac{2}{\sigma^2} \left<y - x, g(x) \right> \\
& \frac{1}{2 \sigma^2} \left(\|g(x)\|^2 - \|g(y)\|^2 \right) \approx
  0. 
\end{align*}
This leads to
\begin{align*}
  \frac{\pi(y)Q(y, x)}{\pi(x) Q(x, y)} \approx \exp \left(\left<y - x, \nabla \log \pi(x) -
    \frac{2g(x)}{\sigma^2} \right> \right) 
\end{align*}
when $y$ is close to $x$. This clearly suggests taking 
\begin{align*}
  g(x) = \frac{\sigma^2}{2} \nabla \log \pi(x). 
\end{align*}
This leads to the MALA proposal:
\begin{align}\label{lang_prop}
  y = x + \frac{1}{2} \sigma^2 \nabla \log \pi(x) + \sigma z ~~ \text{
  where $z \sim N(0, I_d)$}. 
\end{align}
When $\sigma$ is small, $y$ is automatically close to $x$, so the
above argument can be used to show that the Metropolis-Hastings
acceptance ratio for this proposal is nearly equal to 1. 

\section{Mechanics Interpretation}

We can interpret the RWM proposal \eqref{rwm_y} in physical terms
as follows: 
\begin{equation}\label{rwm_phy}
  \begin{split}
  &  x(0) = x, ~~ \dot{x}(0) = z \sim N(0, I_d) \\
  & \text{and, motion with constant velocity $z$ for time $[0, \sigma]$} \\
  & \implies x(\sigma) = \text{RWM proposal}
  \end{split}  
\end{equation}
Similarly, we can interpret the MALA proposal \eqref{lang_prop} in
physical terms as follows: 
\begin{equation}\label{mala_phy}
  \begin{split}
  &  x(0) = x, ~~ \dot{x}(0) = z \sim N(0, I_d) \\
  & \text{and, motion with constant acceleration $\nabla \log \pi(x)$
    for time $[0, \sigma]$} \\ 
  & \implies x(\sigma) = \text{MALA proposal}
  \end{split}  
\end{equation}
The physical analogy makes the difference between RWM and MALA transparent.
In both cases the particle starts at $x$ with a random velocity $z \sim
N(0, I_d)$, and the proposal $y = x(\sigma)$ is the position after time
$\sigma$. The distinction lies in the \emph{dynamics}:
\begin{itemize}
  \item \textbf{RWM} is a free particle: no force acts, so the particle
    travels in a straight line at constant velocity. The proposal is
    purely isotropic and \emph{blind} to the target $\pi$ --- it has no
    tendency to move toward regions of high probability.

  \item \textbf{MALA} subjects the particle to a constant acceleration
    $\nabla \log \pi(x)$, i.e.\ a force pointing in the direction of
    steepest ascent of $\log\pi$. By Newton's second law,
    \[
      x(\sigma) = x + \sigma z + \frac{\sigma^2}{2}\nabla\log\pi(x),
    \]
    which is precisely the Langevin proposal \eqref{lang_prop}. The
    drift term $\frac{\sigma^2}{2}\nabla\log\pi(x)$ biases every
    proposal \emph{uphill} in $\log\pi$, so the chain is steered toward
    high-probability regions before the Metropolis correction is even
    applied.
\end{itemize}
As a consequence, MALA proposals are systematically better aligned with
$\pi$ than RWM proposals. In
practice this allows MALA to take much larger steps $\sigma$ while
maintaining a reasonable acceptance rate, leading to faster exploration
of the target distribution.

In the coming lectures, we shall look at HMC, which is a further
modification of \eqref{mala_phy}. Specifically, the HMC proposal is
obtained via:
\begin{equation}\label{hmc_phy}
  \begin{split}
  &  x(0) = x, ~~ \dot{x}(0) = z \sim N(0, I_d) \\
  & \ddot{x}(t) = \nabla \log \pi(x(t)) \\
  & \implies x(\sigma) = \text{HMC proposal}
  \end{split}  
\end{equation}
The step from MALA to HMC is precisely the step from \emph{constant}
acceleration to \emph{position-dependent} acceleration. In MALA the
acceleration $\nabla\log\pi(x)$ is evaluated once at the current position $x$
and held fixed for the entire flight time $[0,\sigma]$. In HMC the
acceleration is updated continuously along the trajectory: at every time $t$
the particle feels $\nabla\log\pi(x(t))$. One downside to this
time-dependent acceleration is that $x(\sigma)$ usually cannot be
written in closed form as a function of $x$ and $z$ (unlike RWM and
MALA), and some numerical integration scheme has to be used to figure
out $x(\sigma)$. 








%\begin{thebibliography}{99}

%\bibitem[Roberts and Rosenthal(2004)]{RobertsRosenthal2004}
%Roberts, G. O. and Rosenthal, J. S. (2004).
%General state space Markov chains and MCMC algorithms.
%\textit{Probability Surveys}, \textbf{1}, 20--71.

%\bibitem[Meyn and Tweedie(2009)]{MeynTweedie2009}
%Meyn, S. P. and Tweedie, R. L. (2009).
%\textit{Markov Chains and Stochastic Stability}, 2nd ed.
%Cambridge University Press.


  
%\bibitem[Chib and Greenberg(1995)]{chib1995}
%Chib, S. and Greenberg, E. (1995).
%Understanding the Metropolis--Hastings algorithm.
%\textit{The American Statistician}, \textbf{49}, 327--335.


%\bibitem[MacKay and Peto(1995)]{MacKay1995HierarchicalDirichlet}
%D.~J.~C. MacKay and L.~C. Peto,
%\newblock ``A hierarchical Dirichlet language model,''
%\newblock \emph{Natural Language Engineering}, vol.~1,
%no.~3, pp.~289--308, 1995.

%\bibitem[Rubin(1981)]{Rubin1981BayesianBootstrap}
%D.~B. Rubin,
%\newblock ``The Bayesian bootstrap,''
%\newblock \emph{The Annals of Statistics}, vol.~9,
%no.~1, pp.~130--134, 1981.
  
%\bibitem[Fisher(1934)]{Fisher1934}
%R.~A. Fisher,
%\newblock ``Probability, likelihood and quantity of information in the logic of
%uncertain inference,''
%\newblock \emph{Proceedings of the Royal Society of London. Series~A,
%Containing Papers of a Mathematical and Physical Character}, vol.~146,
%no.~856, pp.~1--8, 1934.

%\bibitem[Efron(2010)]{Efron}
%Efron, B. (2010)
%\textit{Large-scale inference: empirical Bayes methods for estimation,
%testing, and prediction}. Cambridge University Press, 2010.

%\bibitem[Shumway and Stoffer(2017)]{ShumwayStoffer}
%Shumway, R. H. and Stoffer, D. S. (2017)
%\textit{Time Series Analysis and Its Applications: With R
%  Examples}. Springer, 4th edition. 
  
%     \bibitem[Jaynes(2003)]{Jaynes} Jaynes, E. T, (2003)
%  \textit{Probability theory: the logic of science}. Cambridge
%  University Press, 2003.
 
%  \bibitem[Sinai(1992)]{Sinai} Sinai, Y. G, (1992)
%  \textit{Probability theory: an introductory course}. Springer
%  Verlag, 1992.  

%\bibitem[{deGroot(1986)}]{DeGrootBlackwell}DeGroot, M. H. (1986)
%       A conversation with David Blackwell.
%\newblock \emph{Statistical Science}, \textbf{1}, 40--53.

%         \bibitem[Mosteller(1987)]{Mosteller} Mosteller, F, (1987)
%  \textit{Fifty challenging problems in probability with
%    solutions}. Courier Corporation, 1987.

%    \bibitem[MacKay(2003)]{MacKay} MacKay, D. J. C., (2003)
%  \textit{Information theory, inference, and learning
%  algorithms}. Cambridge University Press, 2003.

%\bibitem[{Lindley(2013)}]{Lindley}Lindley, D. V. (2013)
%   \textit{Understanding Uncertainty}. John Wiley and sons, 2013.


%\end{thebibliography}







\end{document}

%%% Local Variables:
%%% mode: latex
%%% TeX-master: t
%%% End:
